
\section{Algebra generated by the expectation operators}
\label{sec:6.9}
\begin{proposition}\label{prop:6.5}
The product $(J-N)(M-\tau\omega)$ commutes with $\Omega$.
\end{proposition}
\begin{proof}
Put $r'=1-1/r$ and $\tau'=n^2-\tau$. The identities of
Section~\ref{sec:6.4} give
\[
 \Omega(J-N)(M-\tau\omega)=r'\tau'\Omega.
\]
By symmetry, $(J-N)\pi^{\mathsf T}=r'u^{\mathsf T}$ and
$(M-\tau\omega)u^{\mathsf T}=\tau'\pi^{\mathsf T}$. Hence
\[
 (J-N)(M-\tau\omega)\Omega=r'\tau'\Omega.
\]
\end{proof}

\begin{theorem}\label{thm:6.5}
If $A$ is doubly stochastic, the matrices $M$, $N$ and $J$ generate a
commutative algebra. The same is true of $M-\tau\omega$, $J-N$ and $J$.
\end{theorem}
\begin{proof}
Here $\pi=u/n$, $\Omega=J/n$ and $\omega=I/n$.
The equipartition identity $\pi kN=u/r$ for every kernel element $k$
yields $ukN=nu/r$. Therefore
\[
 MN=\mathbb E[K^{\mathsf T}u^{\mathsf T}uKN]
    =\frac nr\mathbb E[K^{\mathsf T}u^{\mathsf T}]u
    =\frac nr J.
\]
Since $M$ and $N$ are symmetric, transposition gives $NM=MN$.
Also $JM=nJ=MJ$ and $JN=(n/r)J=NJ$.
Subtracting multiples of $I$ and forming $J-N$ preserves commutativity.
\end{proof}

\paragraph{Further question.}
Without double stochasticity, the structure of the commutator $MN-NM$
requires separate analysis. No general commutativity assertion is made here.
